\documentclass[11pt,a4paper]{article}
\usepackage[margin=27mm]{geometry}
\usepackage{amsmath,amssymb,amsthm,booktabs,longtable,hyperref}
\newtheorem{theorem}{Theorem}
\newtheorem{lemma}[theorem]{Lemma}
\newtheorem{corollary}[theorem]{Corollary}
\newcommand{\F}{\mathbb F}
\newcommand{\Sz}{\operatorname{Sz}}
\newcommand{\gen}[1]{\langle #1\rangle}
\title{Maximal subgroups and elementary reductions for Suzuki geometries}
\author{Human Compression project}\date{}
\begin{document}\maketitle
\begin{abstract}
We give complete proofs of the preliminary arithmetic and subgroup
arguments used in the rank-two Suzuki classification. The maximal
subgroups of a point stabilizer follow from its Frattini quotient and
an affine fixed point argument; the two cyclic normalizer cases follow
from the same argument. A general obstruction then excludes a point
stabilizer from every firm residually weakly primitive geometry of rank
greater than one. These are surviving results of the source manuscript,
whose full classification requires a separate correction. Neither the
remaining families nor their global counts are established here.
\end{abstract}

\section{Arithmetic of the three odd factors}
Write $q=2^n$, where $n\ge3$ is odd, and set
\[
 t_q=2^{(n+1)/2},\qquad a_\epsilon(q)=q+\epsilon t_q+1
 \quad(\epsilon\in\{1,-1\}).
\]
Thus $t_q^2=2q$ and
\begin{equation}\label{factor}
             a_+(q)a_-(q)=q^2+1.
\end{equation}

\begin{lemma}\label{arithmetic}
The integers $q-1,a_+(q),a_-(q)$ are pairwise coprime.
If $2^n-1$ is prime, then $n$ is prime.
Suppose $s=2^d\ge8$, with $d$ odd, and $q=s^r$, with $r$ odd.
Then
\[
 s^2+1\mid q^2+1,\qquad
 s^2+1\nmid a_+(q),\qquad s^2+1\nmid a_-(q).
\]
If also $r>1$, all three integers $q-1,a_+(q),a_-(q)$ are composite.
\end{lemma}
\begin{proof}
Every integer in the first assertion is odd. A common divisor of
$a_+(q)$ and $a_-(q)$ divides $2t_q$, a power of two, so is one.
If an odd integer $h$ divides $q-1$ and $a_\epsilon(q)$, reduction
modulo $h$ gives $q=1$, $\epsilon t_q=-2$, and $t_q^2=2$.
Squaring the second equality gives $4=2$, so $h\mid2$ and $h=1$.
If $n=uv$ with $u,v>1$, the identity
\[
 2^{uv}-1=(2^u-1)(1+2^u+\cdots+2^{u(v-1)})
\]
is a nontrivial factorization. This proves the assertion about a
prime $2^n-1$, and proves directly that $q-1$ is composite when
$q=s^r$, $s>2$, $r>1$.

Put $m=s^2+1$ and $t=t_s$. Since $s^2=-1\pmod m$ and $r$ is odd,
$q^2+1=s^{2r}+1$ is zero modulo $m$. Also
$t_q=t s^{(r-1)/2}$. Reducing the odd integer $r$ modulo eight gives
the following representatives for $a_\epsilon(q)$ modulo $m$:
\[
\begin{array}{c|cccc}
 r\pmod8&1&3&5&7\\ \hline
 a_\epsilon(q)&1+s+\epsilon t&1-s+\epsilon st&
                1+s-\epsilon t&1-s-\epsilon st.
\end{array}
\]
All four expressions, for either sign, are nonzero. Indeed,
$s>t>1$, while $st>s-1$. Their absolute values are less than $m$:
$s+t+1<s^2+1$, and $st+s-1<s^2+1$ because $t+1\le s$ for
$s\ge8$. None can therefore represent zero modulo $m$.

Finally assume $r\ge3$. Both $a_\epsilon(q)$ exceed $s^2+1$.
To see this without an asymptotic estimate, $t_q\le q/2$ for
$q\ge8$, so
\[
 a_\epsilon(q)\ge q-t_q+1\ge s^3/2+1>s^2+1.
\]
If one of these two integers were prime, it would be coprime to
$m$, since it exceeds $m$. By \eqref{factor} and $m\mid q^2+1$,
the other factor would then be divisible by $m$, contrary to the
preceding paragraph. Both are composite.
\end{proof}

For completeness, the remaining elementary congruence used in the
source is Fermat's $p\mid2^p-2$ for prime $p$. For odd $p$,
multiplication by $2$ permutes the nonzero residues modulo $p$.
Multiplying those residues and cancelling $(p-1)!$ gives
$2^{p-1}=1\pmod p$. The case $p=2$ is immediate.

Lemma~\ref{arithmetic} proves the mathematical content of source
Lemmas 3.1--3.6, and the last paragraph gives Lemma 3.7. In Lemma
3.2 the proper-subfield condition must be explicit: $r>1$ and $s>2$.
The printed statement does not state $r>1$. Taken literally it fails
at $q=s=8$, $r=1$. This is a correction of its domain, not an
additional refutation of the source's principal classification.
The source attributes Lemmas 3.1 and 3.7 to Cataldi and Fermat,
and the argument for Lemma 3.5 to J.~A.~Thas. We retain those
attributions; the residue table above is a common proof of the same
arithmetic assertions, not a claim of priority.

\section{Two elementary subgroup arguments}
We write $\Phi(P)$ for the intersection of the maximal subgroups of
a finite group $P$. Only the following elementary property of finite
$p$-groups is needed. If $M$ is maximal in a finite $p$-group $P$,
then $M$ is normal and $[P:M]=p$. Here is a proof. A proper subgroup
$K$ acts on the left cosets $P/K$. Its fixed cosets are exactly
$N_P(K)/K$. Every other orbit has size divisible by $p$, and
$|P:K|$ is divisible by $p$. There is at least one fixed coset,
so there are at least $p$, and $N_P(K)>K$. For maximal $M$ this
forces $N_P(M)=P$. The quotient $P/M$ has no proper nontrivial
subgroup, so has order $p$.
Consequently every square in a finite two-group belongs to its
Frattini subgroup. We also use the immediate generating property:
if $J\le P$ and $J\Phi(P)=P$, then $J=P$, since otherwise a
maximal subgroup containing $J$ would contain $J\Phi(P)$.

\begin{lemma}\label{frattini}
Let $P\lhd B$ be a finite $p$-group, and let $Z\le\Phi(P)$ be
normal in $B$. Every maximal subgroup of $B$ contains $Z$.
\end{lemma}
\begin{proof}
If a maximal subgroup $M$ omits $Z$, then $MZ=B$ because $Z$ is
normal. As $Z\le P$, intersecting with $P$ gives
$P=(M\cap P)Z$. Thus $P=(M\cap P)\Phi(P)$, and the generating
property gives $P\le M$. This contradicts $Z\nleq M$.
\end{proof}

\begin{lemma}\label{affine}
Let $V$ be a finite vector space in characteristic $p$, and let
$H\le\operatorname{GL}(V)$ have order prime to $p$. Every complement
to $V$ in $V\rtimes H$ is conjugate by a translation to $H$.
If $H$ acts irreducibly on $V$ over its prime field, the maximal subgroups of $V\rtimes H$
are the inverse images of the maximal subgroups of $H$, together
with the conjugates of $H$.
\end{lemma}
\begin{proof}
A complement $K$ has order $|H|$ and acts affinely on $V$. The point
\[
                 v=|K|^{-1}\sum_{k\in K}k(0)
\]
is fixed by every element of $K$: affine maps preserve affine
averages, and left multiplication permutes $K$. Translating $v$ to
zero conjugates $K$ into $H$; equality follows from their orders.

A maximal subgroup $M$ containing $V$ is the inverse image of a
maximal subgroup of $H$. Otherwise $MV=V\rtimes H$. The additive
subgroup $M\cap V$ is normalized by $M$ and by $V$, hence by the
whole group. It is a vector subspace over the prime field, invariant
under $H$. Irreducibility here is understood over that prime field;
it gives $M\cap V=0$, so $M$ is a complement. Conversely, every
proper overgroup $L$ of $H$ satisfies
$L=(L\cap V)H$. Its intersection with $V$ is invariant, so is either
zero or $V$. Thus $H$ is maximal.
\end{proof}

The prime-field qualification in this lemma matters. We use it for
$V=(\F_q,+)$ and the full scalar group $\F_q^*$, for which it holds:
an additive invariant subgroup containing a nonzero $v$ contains
every scalar multiple of $v$, hence all of $\F_q$.

\begin{corollary}\label{cyclic}
Let $m$ be odd and $d\in\{2,4\}$. In any semidirect product
$C_m\rtimes C_d$, the maximal subgroups are, up to conjugacy,
\[
 C_m\rtimes C_{d/2},\qquad
 C_{m/p}\rtimes C_d\quad(p\text{ a prime divisor of }m).
\]
For $d=2$, the first group is $C_m$. For $m=1$ only the first
group occurs. The statement allows an action that is not faithful.
\end{corollary}
\begin{proof}
Write $A=C_m$. A maximal subgroup containing $A$ is the inverse
image of the unique maximal subgroup of $C_d$. Otherwise $MA$ is
the whole group. Put $K=M\cap A$. Every subgroup of $A$ is
characteristic in $A$, so $K$ is normal in the whole group. In the
quotient by $K$, the image of $M$ is a complement to the cyclic
normal subgroup $A/K$, and is maximal. If $A/K$ had a proper
nontrivial subgroup $T$, that subgroup would be invariant and
$T(M/K)$ would be a proper overgroup of $M/K$, a contradiction.
Thus $|A:K|=p$ is prime. Lemma~\ref{affine}, applied to
$A/K\cong(\F_p,+)$, shows that this complement is conjugate to
the image of the specified $C_d$. Taking inverse images proves
the stated form. Conversely each such subgroup has index $p$,
and so is maximal.
\end{proof}

\section{The point stabilizer}
Let $F=\F_q$, $q=2^n$, $n\ge3$ odd, and put
$\theta=2^{(n+1)/2}$. We use the classical Suzuki root group
\[
 P=\{u(a,b):a,b\in F\},\qquad
 u(a,b)u(c,e)=u(a+c,b+e+ac^\theta).
\]
Let $Z=\{u(0,b):b\in F\}$ and $H=\{h(t):t\in F^*\}$, with
\[
 u(a,b)^{h(t)}=u(t^\theta a,t^{\theta+2}b).
\]
Then the point stabilizer is $B=P\rtimes H$. These formulas and
their realization in the faithful Suzuki matrix group are retained
from the classical construction, also used in Section 2.2 of the
source and in the project's separately reviewed scalar-subspace
note. We do not reconstruct the whole Suzuki group here.

\begin{theorem}\label{borel}
The maximal subgroups of $B$ are, up to conjugacy in $B$,
\[
 P\rtimes C_{(q-1)/p}\quad(p\mid q-1\text{ prime}),
 \qquad Z\rtimes H.
\]
In particular every maximal subgroup of $B$ contains $Z$.
\end{theorem}
\begin{proof}
The multiplication law shows that $Z$ is central in $P$,
$P/Z\cong(F,+)$, and
\[
                       u(a,b)^2=u(0,a^{\theta+1}).
\]
The exponent $\theta+1$ is coprime to $q-1$. Indeed, if an odd
integer $h$ divides both, then $\theta=-1\pmod h$, whereas
$\theta^2=2q=2\pmod h$. It follows that $1=2\pmod h$, so $h=1$.
Every element of $Z$ is consequently a square, and $Z\le\Phi(P)$.
Conversely, the inverse images of the prime-field hyperplanes of
$P/Z$ are maximal subgroups of $P$ whose intersection is $Z$.
Thus $\Phi(P)=Z$.

Lemma~\ref{frattini} shows that every maximal subgroup of $B$
contains $Z$. Modulo $Z$, the conjugation action of $H$ on $F$
is multiplication by $t^\theta$. This ranges over all nonzero
scalars, so is irreducible over $\F_2$. Lemma~\ref{affine}
classifies the maximal subgroups of $B/Z$. Their inverse images
are exactly the groups displayed above.
\end{proof}

Theorem~\ref{borel} gives source Lemma 5.1, including the
completeness of its list. Corollary~\ref{cyclic} gives Lemma 5.2
for the cyclic normalizers of order $4a_\epsilon(q)$, and Lemma
5.3 for the dihedral group $C_{q-1}\rtimes C_2$. Thus the three
lists use one affine complement argument. Their assertions remain
attributed to the source; this note supplies the complete common
derivation and does not claim that these elementary tools are new.

\section{Excluding a point stabilizer in every rank}
We first isolate a general geometric obstruction. In a flag-transitive
coset geometry
\[\Gamma(G;(G_i)_{i\in I}),\]
the residue of the base element of type $0$ has acting group $G_0$ and stabilizers
$G_0\cap G_i$, $i\ne0$. A transitive coset action is primitive
if and only if its point stabilizer is maximal: blocks containing
the base coset correspond to intermediate subgroups.

\begin{lemma}\label{terminal}
Let $G$ be finite and $B<G$ be maximal. Suppose that every maximal subgroup $J$ of
$B$ is contained in no maximal subgroup of $G$ other than $B$.
Then no firm flag-transitive residually weakly primitive geometry
of rank greater than one has $B$ as one of its parabolic subgroups.
\end{lemma}
\begin{proof}
Suppose $G_0=B$. Weak primitivity of the residue at the base
element of type $0$ gives an $i\ne0$ such that
$J=B\cap G_i$ is maximal in $B$. Every proper subgroup of the
finite group $G$ containing $J$ lies in a maximal subgroup, hence
lies in $B$. Thus $G_i\le B$, so $G_i=J$.
Set $G_{I\setminus\{0\}}=\bigcap_{j\ne0}G_j$ and
$G_I=\bigcap_jG_j$. Since $G_i\le B$, these two intersections
are equal. The base flag of all types except $0$ therefore extends
to exactly
\[
             [G_{I\setminus\{0\}}:G_I]=1
\]
chamber. This contradicts firmness. The same proof applies when
the geometry has rank two; no truncation assumption is needed.
\end{proof}

\begin{theorem}\label{exclude}
For $G=\Sz(q)$, $q=2^n$, $n\ge3$ odd, the point stabilizer $B$
cannot occur as a parabolic subgroup in a firm flag-transitive
residually weakly primitive geometry of rank greater than one.
\end{theorem}
\begin{proof}
We use Suzuki's classical list of maximal subgroups, quoted as
Table 1 of the source: point stabilizers, $D_{2(q-1)}$,
$C_{a_+(q)}\rtimes C_4$, $C_{a_-(q)}\rtimes C_4$, and maximal
simple subfield groups $\Sz(s)$ with $q=s^r$, $r$ an odd prime.
Let $J$ be maximal in $B$. Theorem~\ref{borel} gives $Z\le J$,
where $Z$ is elementary abelian of order $q\ge8$.

The dihedral maximal subgroup has Sylow two-subgroups of order
two. The two cyclic normalizers have cyclic Sylow two-subgroups
of order four. None can contain $Z$. A simple proper subfield
group has Sylow two-subgroups of order $s^2$, whereas
$q=s^r\ge s^3>s^2$. It too cannot contain $Z$.

Finally, a nonidentity element of the root group $P$ fixes exactly
one point in the classical ovoid action. In the affine coordinates
on the complement of the distinguished point, its action is
\[
                  (x,y)\longmapsto(x+a,y+b+a^\theta x).
\]
A fixed affine point forces $a=0$, then $b=0$, so no nonidentity
element fixes such a point. Consequently every point stabilizer
containing a fixed nonidentity element of $Z$ is $B$ itself.
The only maximal subgroup of $G$ containing $J$ is therefore
$B$, and Lemma~\ref{terminal} applies.
\end{proof}

This proves the full rank range of source Lemma 6.1. The geometric
step is written out because merely excluding rank-two configurations
does not, by itself, justify an assertion in every rank.

\section{Attribution and the remaining source obligations}
The source is D.~Leemans, \emph{The Rank 2 Geometries of the Simple
Suzuki Groups $\Sz(q)$}, Beitr\"age zur Algebra und Geometrie
39(1) (1998), 97--120, in the fixed twenty-four-page edition.
The assertions treated here are its seven arithmetic lemmas,
three maximal-subgroup lemmas, and Lemma 6.1. We retain Suzuki's
maximal-subgroup theorem and the standard faithful root-group and
ovoid constructions as explicit premises. The source credits
Suzuki's \emph{On a class of doubly transitive groups}, Annals of
Mathematics 75 (1962), 105--145, and the treatments by Tits and
L\"uneburg for this structure. No fresh proof of that classification
or new historical priority is asserted.

The common proofs here use elementary finite group theory:
the Frattini generating property, irreducibility of the full scalar
action, affine averaging, and the correspondence between blocks
and intermediate subgroups. The distinct scalar-subspace correction
to the source's mixed-parabolic family remains unchanged. The
remaining assertions in Section 6, the corrected global family
list, Tables 2--12 and their counts still require their own
source review. This component supplies neither a complete
treatment of the original nor an open-conjecture resolution.
No Lean formalization or compression ratio is claimed for it.
\end{document}
